% 第31批50道高难综合题：当前为审核候选，未完成独立专家严格质检。

% 每题列出四个关键推导结果与物理检验线索；正式发布前仍需详细逐问复算。

\subsection*{M202\quad 中心势小振动与轨道频率}

\textbf{模型与问题：} 有效势$U_{eff}=L^2/(2mr^2)+ar^s$，$a s>0$且存在稳定圆轨道。 (1)建立控制关系；(2)求核心解；(3)分析相关物理量；(4)讨论守恒、极限或稳定性。

\par\textbf{推导要点：} (1) 令$U'_{eff}=0$得$L^2/m=asr_c^{s+2}$； (2) $\Omega^2=L^2/(m^2r_c^4)=asr_c^{s-2}/m$； (3) $\kappa^2=U''_{eff}(r_c)/m=(s+2)\Omega^2$； (4) 稳定要求$s+2>0$，径向周期的进动$\Delta\phi=2\pi\Omega/\kappa$。。

\par\medskip

\subsection*{M203\quad 张紧弦转动端点稳定}

\textbf{模型与问题：} 半径$r$的质量$m$环绕中心以$L$角动量运动，附径向弹簧势$k(r-r_0)^2/2$。 (1)建立控制关系；(2)求核心解；(3)分析相关物理量；(4)讨论守恒、极限或稳定性。

\par\textbf{推导要点：} (1) $U_{eff}=L^2/(2mr^2)+k(r-r_0)^2/2$； (2) 平衡$r_*$满足$kr_*^3(r_*-r_0)=L^2/m$； (3) 径向频率$\omega_r^2=[3L^2/(mr_*^4)+k]/m$； (4) 角频率$\Omega=L/(mr_*^2)$，均由有效势二阶展开得出。。

\par\medskip

\subsection*{M204\quad 各向异性二维振子守恒}

\textbf{模型与问题：} 势$U=m(\omega_x^2x^2+\omega_y^2y^2)/2$且频率之比为有理数$p/q$。 (1)建立控制关系；(2)求核心解；(3)分析相关物理量；(4)讨论守恒、极限或稳定性。

\par\textbf{推导要点：} (1) 正则方程分别为$\ddot x+\omega_x^2x=0$与$\ddot y+\omega_y^2y=0$； (2) 运动为$A\cos(\omega_xt+\alpha)$及$B\cos(\omega_yt+\beta)$； (3) 轨道闭合的必要充分条件为$\omega_x/\omega_y\in\mathbb Q$（两振幅非零）； (4) 最小共同周期$T=2\pi p/\omega_x=2\pi q/\omega_y$若$\omega_x/\omega_y=p/q$，应改为$T=2\pi p/\omega_x=2\pi q/\omega_y$。。

\par\medskip

\subsection*{M205\quad 旋转参考系的惯性波}

\textbf{模型与问题：} 不可压无黏液体以$\boldsymbol\Omega=\Omega\hat z$匀速旋转，线性速度扰动为$\mathbf u e^{i(\mathbf k\cdot\mathbf x-\omega t)}$。 (1)建立控制关系；(2)求核心解；(3)分析相关物理量；(4)讨论守恒、极限或稳定性。

\par\textbf{推导要点：} (1) 线性动量方程$-i\omega\mathbf u+2\boldsymbol\Omega\times\mathbf u=-i\mathbf k p/\rho$； (2) 不可压条件$\mathbf k\cdot\mathbf u=0$； (3) 取与$\mathbf k$正交投影得惯性波$\omega^2=4\Omega^2k_z^2/k^2$； (4) 故$|\omega|\le2\Omega$，当$k_z=0$频率退化为零。。

\par\medskip

\subsection*{M206\quad 耦合振子的正则模}

\textbf{模型与问题：} 相同质量$m$的两振子各与固定墙由弹簧$k$相连，两质点间弹簧$k_c$。 (1)建立控制关系；(2)求核心解；(3)分析相关物理量；(4)讨论守恒、极限或稳定性。

\par\textbf{推导要点：} (1) 方程$m\ddot x_1=-(k+k_c)x_1+k_cx_2$及对称式； (2) 对称模$(1,1)$频率$\omega_+=\sqrt{k/m}$； (3) 反对称模$(1,-1)$频率$\omega_-=\sqrt{(k+2k_c)/m}$； (4) 一般初态分解为两正交模的线性组合且总能量保持。。

\par\medskip

\subsection*{M207\quad 刚体无力矩 Euler 旋转}

\textbf{模型与问题：} 轴对称刚体$I_1=I_2=I_\perp$、$I_3$，无外力矩。 (1)建立控制关系；(2)求核心解；(3)分析相关物理量；(4)讨论守恒、极限或稳定性。

\par\textbf{推导要点：} (1) Euler方程$I_\perp\dot\omega_1=(I_\perp-I_3)\omega_2\omega_3$及$I_\perp\dot\omega_2=(I_3-I_\perp)\omega_3\omega_1$； (2) $\dot\omega_3=0$； (3) 横向分量以$|I_3-I_\perp||\omega_3|/I_\perp$匀速进动； (4) 能量$\sum I_i\omega_i^2/2$与角动量模守恒。。

\par\medskip

\subsection*{M208\quad 重力梯度势的潮汐稳定}

\textbf{模型与问题：} 两点质量以小位移$\boldsymbol\xi$围绕半径$r$的中心引力轨道，中心质量$M$。 (1)建立控制关系；(2)求核心解；(3)分析相关物理量；(4)讨论守恒、极限或稳定性。

\par\textbf{推导要点：} (1) 引力场$\mathbf g=-GM\mathbf r/r^3$； (2) 线性化$\delta\mathbf g=-GM[\boldsymbol\xi-3(\hat r\cdot\boldsymbol\xi)\hat r]/r^3$； (3) 径向相对加速度$+2GM\xi_r/r^3$而横向为$-GM\boldsymbol\xi_\perp/r^3$； (4) 这分别对应径向拉伸和横向压缩，忽略轨道系科里奥利项。。

\par\medskip

\subsection*{M209\quad 含速度阻尼的受迫振动共振}

\textbf{模型与问题：} 振子$m\ddot x+b\dot x+kx=F_0\cos\omega t$，$b,k,m>0$。 (1)建立控制关系；(2)求核心解；(3)分析相关物理量；(4)讨论守恒、极限或稳定性。

\par\textbf{推导要点：} (1) 复幅$X=F_0/(k-m\omega^2+ib\omega)$（以$e^{i\omega t}$约定）； (2) 幅值$A=F_0/\sqrt{(k-m\omega^2)^2+b^2\omega^2}$； (3) 峰值在$\omega_{pk}^2=k/m-b^2/(2m^2)>0$时存在； (4) 平均耗散$\bar P=b\omega^2A^2/2$，与输入平均功率相等。。

\par\medskip

\subsection*{M210\quad 广义坐标的参数共振}

\textbf{模型与问题：} 弱参数激励$\ddot x+\omega_0^2[1+h\cos(2\omega_0t)]x=0$，$|h|\ll1$。 (1)建立控制关系；(2)求核心解；(3)分析相关物理量；(4)讨论守恒、极限或稳定性。

\par\textbf{推导要点：} (1) 设$x=A(t)\cos\omega_0t+B(t)\sin\omega_0t$并忽略二阶慢导数； (2) 平均后$\dot A=-(h\omega_0/4)B$、$\dot B=-(h\omega_0/4)A$（振幅基底符号可整体约定）； (3) 特征增长率$|h|\omega_0/4$； (4) 离共振失谐使增长受限，非线性及耗散会限制实际增长。。

\par\medskip

\subsection*{E202\quad 有限半径线圈轴场曲率}

\textbf{模型与问题：} 半径$R$电流$I$的圆环在$z$轴上产生磁场。 (1)建立控制关系；(2)求核心解；(3)分析相关物理量；(4)讨论守恒、极限或稳定性。

\par\textbf{推导要点：} (1) Biot–Savart积分给$B_z=\mu_0IR^2/[2(R^2+z^2)^{3/2}]$； (2) 中心场$B_0=\mu_0I/(2R)$； (3) 近中心展开$B_z/B_0=1-3z^2/(2R^2)+O(z^4)$； (4) 梯度在$z=0$为零，二阶曲率为$-3B_0/R^2$。。

\par\medskip

\subsection*{E203\quad 磁偶极子场的多极结构}

\textbf{模型与问题：} 点磁偶极矩$\mathbf m$处在真空，观察点$r>0$。 (1)建立控制关系；(2)求核心解；(3)分析相关物理量；(4)讨论守恒、极限或稳定性。

\par\textbf{推导要点：} (1) 矢势$\mathbf A=\mu_0(\mathbf m\times\mathbf r)/(4\pi r^3)$； (2) 旋度得$\mathbf B=\mu_0[3(\mathbf m\cdot\hat r)\hat r-\mathbf m]/(4\pi r^3)$； (3) 轴上场大小$\mu_0(2m)/(4\pi r^3)$； (4) 赤道面场方向反于$\mathbf m$，大小$\mu_0m/(4\pi r^3)$。。

\par\medskip

\subsection*{E204\quad 球形导体的电势镜像法}

\textbf{模型与问题：} 半径$a$接地金属球外点电荷$q$位于$r_0>a$轴上。 (1)建立控制关系；(2)求核心解；(3)分析相关物理量；(4)讨论守恒、极限或稳定性。

\par\textbf{推导要点：} (1) 在同轴$r'=a^2/r_0$放像电荷$q'=-qa/r_0$； (2) 两电荷势在$r=a$处因距离比相等而相消； (3) 真实区域$r\ge a$势为$\varphi=(q/|\mathbf r-r_0\hat z|+q'/|\mathbf r-r'\hat z|)/(4\pi\epsilon_0)$； (4) 总感应电荷由远场单极项给$Q_{ind}=q'$。。

\par\medskip

\subsection*{E205\quad 非线性介质能量密度}

\textbf{模型与问题：} 静电介质本构$\mathbf D=(\epsilon+\chi E^2)\mathbf E$，$\epsilon>0,\chi\ge0$，局域各向同性。 (1)建立控制关系；(2)求核心解；(3)分析相关物理量；(4)讨论守恒、极限或稳定性。

\par\textbf{推导要点：} (1) 准静态可逆充电有$du=\mathbf E\cdot d\mathbf D$； (2) $dD/dE=\epsilon+3\chi E^2$； (3) 积分$u(E)=\epsilon E^2/2+3\chi E^4/4$； (4) 须区分其与$\int_0^E D(E')dE'=\epsilon E^2/2+\chi E^4/4$，后者对应不同热力学势。。

\par\medskip

\subsection*{E206\quad RLC受迫品质因数}

\textbf{模型与问题：} 串联$R,L,C$受驱动电压$V_0\cos\omega t$。 (1)建立控制关系；(2)求核心解；(3)分析相关物理量；(4)讨论守恒、极限或稳定性。

\par\textbf{推导要点：} (1) 阻抗$Z=R+i(\omega L-1/(\omega C))$； (2) 电流幅值$I_0=V_0/|Z|$； (3) 电流共振$\omega_0=1/\sqrt{LC}$，此时平均耗散$V_0^2/(2R)$； (4) 高品质因数下$Q=\omega_0L/R$、半功率带宽$\Delta\omega=R/L$。。

\par\medskip

\subsection*{E207\quad 运动导体棒的能量平衡}

\textbf{模型与问题：} 长度$\ell$的棒在均匀磁场$B$中沿导轨速率$v$运动，回路电阻$R$。 (1)建立控制关系；(2)求核心解；(3)分析相关物理量；(4)讨论守恒、极限或稳定性。

\par\textbf{推导要点：} (1) 磁通变化率给电动势$\mathcal E=B\ell v$； (2) 电流$I=B\ell v/R$； (3) 磁阻力$F=B\ell I=B^2\ell^2v/R$反向速度； (4) 机械功率$Fv=B^2\ell^2v^2/R=I^2R$，正好等于焦耳热。。

\par\medskip

\subsection*{E208\quad 等离子体横向磁波}

\textbf{模型与问题：} 均匀电中性理想MHD流体密度$\rho$、背景磁场$B_0\hat z$，不可压横向扰动沿$z$传播。 (1)建立控制关系；(2)求核心解；(3)分析相关物理量；(4)讨论守恒、极限或稳定性。

\par\textbf{推导要点：} (1) 动量线性化$\rho\partial_t\mathbf u_\perp=(B_0/\mu_0)\partial_z\mathbf b_\perp$； (2) 感应方程$\partial_t\mathbf b_\perp=B_0\partial_z\mathbf u_\perp$； (3) 消元得$\partial_t^2\mathbf u_\perp=v_A^2\partial_z^2\mathbf u_\perp$，$v_A=B_0/\sqrt{\mu_0\rho}$； (4) 波动能的动能和磁能均分。。

\par\medskip

\subsection*{E209\quad 位移电流与电容充电}

\textbf{模型与问题：} 平行板面积$A$的理想电容间距$d$，介电常数$\epsilon$，忽略边缘效应，充电电流$I(t)$。 (1)建立控制关系；(2)求核心解；(3)分析相关物理量；(4)讨论守恒、极限或稳定性。

\par\textbf{推导要点：} (1) 自由面电荷$Q(t)$满足$\dot Q=I$； (2) 板间位移场$D=Q/A$； (3) 位移电流$\int\partial_tD\,dA=I$，与外接导线传导电流连续； (4) 磁场Ampere–Maxwell环路积分对跨越导线或板间的曲面给一致结果。。

\par\medskip

\subsection*{E210\quad 磁场中载流子动力学张量}

\textbf{模型与问题：} 单一带电粒子载流子密度$n$，电荷$q$、质量$m$，弛豫时间$\tau$，均匀$B\hat z$。 (1)建立控制关系；(2)求核心解；(3)分析相关物理量；(4)讨论守恒、极限或稳定性。

\par\textbf{推导要点：} (1) 稳态Drude方程$\mathbf v/\tau-(q/m)\mathbf v\times\mathbf B=(q/m)\mathbf E$； (2) 令$\beta=qB\tau/m$，则$v_x=(q\tau/m)(E_x+\beta E_y)/(1+\beta^2)$、$v_y=(q\tau/m)(E_y-\beta E_x)/(1+\beta^2)$（对交叉乘积约定检验符号）； (3) 电流$\mathbf J=nq\mathbf v$； (4) $\beta\to0$时恢复各向同性$\sigma_0=nq^2\tau/m$。。

\par\medskip

\subsection*{T202\quad 有限热容热机最大可用功}

\textbf{模型与问题：} 两个热容$C_1,C_2$的物体初温$T_1>T_2$，与环境绝热，热机循环运行至同温。 (1)建立控制关系；(2)求核心解；(3)分析相关物理量；(4)讨论守恒、极限或稳定性。

\par\textbf{推导要点：} (1) 可逆条件$C_1\ln(T_f/T_1)+C_2\ln(T_f/T_2)=0$； (2) 解得$T_f=T_1^{C_1/C}T_2^{C_2/C}$、$C=C_1+C_2$； (3) 输出功$W=C_1T_1+C_2T_2-CT_f$； (4) AM-GM证明$W\ge0$，不可逆接触的热机可获功小于此上界。。

\par\medskip

\subsection*{T203\quad van der Waals 临界点}

\textbf{模型与问题：} 一摩尔流体满足$(P+a/V^2)(V-b)=RT$，$a,b>0$。 (1)建立控制关系；(2)求核心解；(3)分析相关物理量；(4)讨论守恒、极限或稳定性。

\par\textbf{推导要点：} (1) 写$P=RT/(V-b)-a/V^2$； (2) 临界条件$(\partial P/\partial V)_T=(\partial^2P/\partial V^2)_T=0$； (3) 联立得$V_c=3b$、$T_c=8a/(27Rb)$、$P_c=a/(27b^2)$； (4) 压缩因子$P_cV_c/(RT_c)=3/8$。。

\par\medskip

\subsection*{T204\quad 热容跃迁与自由能连续性}

\textbf{模型与问题：} 两相在连续相变点$T_c$处$F_1(T_c)=F_2(T_c)$且熵连续。 (1)建立控制关系；(2)求核心解；(3)分析相关物理量；(4)讨论守恒、极限或稳定性。

\par\textbf{推导要点：} (1) 熵$S=-\partial F/\partial T$，故一阶导也连续； (2) 定容热容$C_V=-T\partial^2F/\partial T^2$； (3) 若热容有有限跳跃，则$F$的二阶导有跳跃； (4) 因潜热$L=T_c\Delta S=0$，属于非一阶相变（需另核验热力学极限）。。

\par\medskip

\subsection*{T205\quad 等温压缩的最小外功}

\textbf{模型与问题：} 理想气体$n$摩尔在热库$T$下从$V_1$压至$V_2<V_1$。 (1)建立控制关系；(2)求核心解；(3)分析相关物理量；(4)讨论守恒、极限或稳定性。

\par\textbf{推导要点：} (1) 可逆功$W_{by}=nRT\ln(V_2/V_1)$； (2) 理想气体内能不变，$Q=W_{by}$； (3) 熵变$\Delta S=nR\ln(V_2/V_1)$； (4) 最小外功$-W_{by}=nRT\ln(V_1/V_2)$，不可逆压缩所需功通常更大。。

\par\medskip

\subsection*{T206\quad 光子气体绝热标度}

\textbf{模型与问题：} 腔中处于热平衡的黑体辐射，体积$V$、温度$T$，能量$U=aVT^4$。 (1)建立控制关系；(2)求核心解；(3)分析相关物理量；(4)讨论守恒、极限或稳定性。

\par\textbf{推导要点：} (1) 辐射压力$P=U/(3V)=aT^4/3$； (2) 第一定律可逆绝热$dU+PdV=0$； (3) 代入得$4aVT^3dT+(4/3)aT^4dV=0$； (4) 积分$TV^{1/3}=\mathrm{const}$，因而$PV^{4/3}=\mathrm{const}$。。

\par\medskip

\subsection*{T207\quad 弱温差可逆可用能}

\textbf{模型与问题：} 两恒热容相同$C$的物体初温$T\pm\delta$，$0<\delta\ll T$。 (1)建立控制关系；(2)求核心解；(3)分析相关物理量；(4)讨论守恒、极限或稳定性。

\par\textbf{推导要点：} (1) 熵守恒得$T_f=\sqrt{T^2-\delta^2}$； (2) 可逆取功$W=2C(T-T_f)$； (3) 二项式展开$W=C\delta^2/T+O(C\delta^4/T^3)$； (4) 两者直接接触终温为$T$且总熵增$\Delta S=C\ln[T^2/(T^2-\delta^2)]>0$。。

\par\medskip

\subsection*{T208\quad 不可逆热传导的熵产}

\textbf{模型与问题：} 一维均匀介质$T(x)>0$，导热系数$k>0$，局域平衡。 (1)建立控制关系；(2)求核心解；(3)分析相关物理量；(4)讨论守恒、极限或稳定性。

\par\textbf{推导要点：} (1) Fourier定律$q_x=-k\partial_xT$； (2) 局域熵产密度$\sigma_s=q_x\partial_x(1/T)$； (3) 故$\sigma_s=k(\partial_xT)^2/T^2\ge0$； (4) 均匀温度梯度消失则熵产为零，单位为$\mathrm{W}/(\mathrm{m}^3\mathrm K)$。。

\par\medskip

\subsection*{T209\quad 热涨落的能量方差}

\textbf{模型与问题：} 正则系综配分函数$Z(\beta)$，$\beta=1/(k_BT)$。 (1)建立控制关系；(2)求核心解；(3)分析相关物理量；(4)讨论守恒、极限或稳定性。

\par\textbf{推导要点：} (1) 平均能量$U=-\partial_\beta\ln Z$； (2) 能量方差$\operatorname{Var}(E)=\partial_\beta^2\ln Z$； (3) 由$\partial_\beta=-k_BT^2\partial_T$得到$\operatorname{Var}(E)=k_BT^2C_V$； (4) 正则系综$C_V\ge0$，负热容须讨论非正则体系。。

\par\medskip

\subsection*{W202\quad 晶格边界缺陷局域模}

\textbf{模型与问题：} 一维等质量$m$原子链弹簧$K$，孤立缺陷质量$m_d<m$。 (1)建立控制关系；(2)求核心解；(3)分析相关物理量；(4)讨论守恒、极限或稳定性。

\par\textbf{推导要点：} (1) 体内方程$m\ddot u_n=K(u_{n+1}+u_{n-1}-2u_n)$； (2) 高于声学带边的局域模可设$u_n=A(-1)^ne^{-\kappa|n|}$； (3) 体方程给$m\omega^2=2K(1+\cosh\kappa)$； (4) 缺陷处$m_d\omega^2=2K(1+e^{-\kappa})$，联立确定$\kappa>0$的条件。。

\par\medskip

\subsection*{W203\quad 两支耦合波导的拍长}

\textbf{模型与问题：} 理想两相同波导的模振幅满足$i\,d\mathbf a/dz=\kappa\sigma_x\mathbf a$，$\kappa>0$。 (1)建立控制关系；(2)求核心解；(3)分析相关物理量；(4)讨论守恒、极限或稳定性。

\par\textbf{推导要点：} (1) 矩阵平方$(\sigma_x)^2=I$； (2) 传播算符$U(z)=\cos\kappa z\,I-i\sin\kappa z\,\sigma_x$； (3) 初始$(1,0)^T$时$P_2=\sin^2\kappa z$； (4) 全交换距离$L_c=\pi/(2\kappa)$，总功率守恒。。

\par\medskip

\subsection*{W204\quad 连续介质弹性剪切波}

\textbf{模型与问题：} 均匀各向同性弹性体密度$\rho$、Lamé系数$\lambda,\mu$。 (1)建立控制关系；(2)求核心解；(3)分析相关物理量；(4)讨论守恒、极限或稳定性。

\par\textbf{推导要点：} (1) Navier方程$\rho\ddot{\mathbf u}=(\lambda+\mu)\nabla(\nabla\cdot\mathbf u)+\mu\nabla^2\mathbf u$； (2) 纵波速度$\sqrt{(\lambda+2\mu)/\rho}$； (3) 横波速度$\sqrt{\mu/\rho}$； (4) 稳定弹性要求$\mu>0$及体积模量$\lambda+2\mu/3>0$。。

\par\medskip

\subsection*{W205\quad 非线性波包自相位调制}

\textbf{模型与问题：} 无色散的Kerr波导满足$\partial_zA=i\gamma|A|^2A$，$\gamma$实数。 (1)建立控制关系；(2)求核心解；(3)分析相关物理量；(4)讨论守恒、极限或稳定性。

\par\textbf{推导要点：} (1) 对$|A|^2$求导得$\partial_z|A|^2=0$； (2) 因此$A(z,t)=A(0,t)e^{i\gamma|A(0,t)|^2z}$； (3) 非线性相位$\phi_{NL}=\gamma|A(0,t)|^2z$； (4) 瞬时频率偏移$\delta\omega=-\partial_t\phi_{NL}=-\gamma z\partial_t|A(0,t)|^2$。。

\par\medskip

\subsection*{W206\quad 黏滞声波的复波数}

\textbf{模型与问题：} 一维小振幅声压场满足$\partial_t^2p-c^2\partial_x^2p=\Gamma\partial_t\partial_x^2p$，$\Gamma>0$。 (1)建立控制关系；(2)求核心解；(3)分析相关物理量；(4)讨论守恒、极限或稳定性。

\par\textbf{推导要点：} (1) 代入$e^{i(kx-\omega t)}$得$\omega^2=(c^2-i\Gamma\omega)k^2$； (2) 实频驱动下$k=\omega/\sqrt{c^2-i\Gamma\omega}$； (3) 小损耗展开$k\simeq\omega/c+i\Gamma\omega^2/(2c^3)$； (4) 因$e^{ikx}$有$e^{-\mathrm{Im}k\,x}$，传播向右时振幅衰减。。

\par\medskip

\subsection*{W207\quad 慢变介质的WKB波幅}

\textbf{模型与问题：} 一维无耗散方程$y_{tt}=\partial_x[c(x)^2y_x]$，介质缓变且驱动频率$\omega$。 (1)建立控制关系；(2)求核心解；(3)分析相关物理量；(4)讨论守恒、极限或稳定性。

\par\textbf{推导要点：} (1) 取$y=\Re[A(x)e^{i\int kdx-i\omega t}]$； (2) 主阶色散$k=\omega/c(x)$； (3) 下一阶通量守恒给$A^2c^2k=\mathrm{const}$； (4) 故$A\propto c^{-1/2}$，仅在$|k'|\ll k^2$的WKB区成立。。

\par\medskip

\subsection*{W208\quad 自由表面重力波色散}

\textbf{模型与问题：} 深度$h$、密度$\rho$的不可压无旋流体自由表面小振幅波，重力$g$，忽略表面张力。 (1)建立控制关系；(2)求核心解；(3)分析相关物理量；(4)讨论守恒、极限或稳定性。

\par\textbf{推导要点：} (1) 速度势满足Laplace方程并在底部满足$\phi_z(-h)=0$； (2) 线性自由面动力边界$\phi_t+g\eta=0$； (3) 运动学边界$\eta_t=\phi_z(0)$； (4) 消元得$\omega^2=gk\tanh(kh)$，深水为$gk$、浅水为$ghk^2$。。

\par\medskip

\subsection*{W209\quad 群速度与色散能流}

\textbf{模型与问题：} 均匀无耗散弹性介质中谐波模具有色散$\omega(k)$及每周期平均能密度$\bar u$。 (1)建立控制关系；(2)求核心解；(3)分析相关物理量；(4)讨论守恒、极限或稳定性。

\par\textbf{推导要点：} (1) 窄带波包的包络满足$x-v_gt=\mathrm{const}$，$v_g=d\omega/dk$； (2) 在无耗散、正定能量和合适归一化下平均能流$\bar S=v_g\bar u$； (3) 相速$v_p=\omega/k$一般不同于群速； (4) 近带边$v_g\to0$使净能量传输趋零，结论需无吸收条件。。

\par\medskip

\subsection*{O202\quad 四分之一波长增透膜}

\textbf{模型与问题：} 空气$n_0$、膜$n_1$、基底$n_2$且均为无耗散非磁性，正入射。 (1)建立控制关系；(2)求核心解；(3)分析相关物理量；(4)讨论守恒、极限或稳定性。

\par\textbf{推导要点：} (1) 两界面反射幅系数$r_{01}=(n_0-n_1)/(n_0+n_1)$、$r_{12}=(n_1-n_2)/(n_1+n_2)$； (2) 膜相位$\delta=2\pi n_1d/\lambda$； (3) 总反射$r=(r_{01}+r_{12}e^{2i\delta})/(1+r_{01}r_{12}e^{2i\delta})$； (4) 若$n_1^2=n_0n_2$且$d=\lambda/(4n_1)$，则$r=0$。。

\par\medskip

\subsection*{O203\quad 法布里珀罗腔线宽}

\textbf{模型与问题：} 两相同无耗散镜面强度反射率$R<1$的Fabry–Pérot腔，单次往返相移$\delta$。 (1)建立控制关系；(2)求核心解；(3)分析相关物理量；(4)讨论守恒、极限或稳定性。

\par\textbf{推导要点：} (1) 多束透射几何级数求得透射率$\mathcal T=1/[1+F\sin^2(\delta/2)]$，$F=4R/(1-R)^2$； (2) 峰位$\delta=2\pi m$； (3) 高反射率下相位半高全宽约$4/\sqrt F$； (4) 自由光谱范围与线宽比约$\mathcal F\simeq\pi\sqrt R/(1-R)$。。

\par\medskip

\subsection*{O204\quad 高斯光束瑞利长度}

\textbf{模型与问题：} 真空中基模高斯光束波长$\lambda$，腰半径$w_0$。 (1)建立控制关系；(2)求核心解；(3)分析相关物理量；(4)讨论守恒、极限或稳定性。

\par\textbf{推导要点：} (1) 波矢$k=2\pi/\lambda$，瑞利长度$z_R=\pi w_0^2/\lambda$； (2) 束宽$w(z)=w_0\sqrt{1+(z/z_R)^2}$； (3) 波前曲率$R(z)=z[1+(z_R/z)^2]$； (4) 远场半角发散$\theta\simeq\lambda/(\pi w_0)$，体现腰径与发散反比。。

\par\medskip

\subsection*{O205\quad 偏振态 Jones 计算}

\textbf{模型与问题：} 四分之一波片快轴沿$x$、慢轴沿$y$，入射Jones矢量$(1,1)^T/\sqrt2$。 (1)建立控制关系；(2)求核心解；(3)分析相关物理量；(4)讨论守恒、极限或稳定性。

\par\textbf{推导要点：} (1) 忽略整体相位时波片$J=\mathrm{diag}(1,i)$； (2) 出射场$(1,i)^T/\sqrt2$； (3) 两正交分量等幅、相差$\pi/2$，为圆偏振； (4) 其旋向符号依时间因子$e^{-i\omega t}$及观察方向约定。。

\par\medskip

\subsection*{O206\quad 空间频率与光学傅里叶衍射}

\textbf{模型与问题：} 夫琅禾费近似下入射孔径场$U(x)$。 (1)建立控制关系；(2)求核心解；(3)分析相关物理量；(4)讨论守恒、极限或稳定性。

\par\textbf{推导要点：} (1) 远场复振幅$\widetilde U(k_x)=\int U(x)e^{-ik_xx}dx$； (2) 观察角满足$k_x=(2\pi/\lambda)\sin\theta$； (3) 矩形孔宽$a$给$\widetilde U\propto a\operatorname{sinc}(k_xa/2)$； (4) 首暗纹满足$a\sin\theta=\lambda$，宽度随口径反比。。

\par\medskip

\subsection*{O207\quad 脉冲色散的群时延}

\textbf{模型与问题：} 长度$L$介质折射率$n(\omega)$，相位$\phi(\omega)=n(\omega)\omega L/c$。 (1)建立控制关系；(2)求核心解；(3)分析相关物理量；(4)讨论守恒、极限或稳定性。

\par\textbf{推导要点：} (1) 群时延$\tau_g=d\phi/d\omega=L[n+\omega n']/c$； (2) 群速度$v_g=c/(n+\omega n')$； (3) 群延迟色散$\phi''=L(2n'+\omega n'')/c$； (4) 窄带脉冲近似在大$|\phi''|$时展宽，异常群速并不直接确定信号前沿速度。。

\par\medskip

\subsection*{O208\quad 部分相干双缝可见度}

\textbf{模型与问题：} 双缝平均强度$I_1,I_2$，复相干度$\gamma$。 (1)建立控制关系；(2)求核心解；(3)分析相关物理量；(4)讨论守恒、极限或稳定性。

\par\textbf{推导要点：} (1) 叠加强度$I=I_1+I_2+2\sqrt{I_1I_2}\Re(\gamma e^{i\phi})$； (2) 极值$I_{max/min}=I_1+I_2\pm2\sqrt{I_1I_2}|\gamma|$； (3) 可见度$V=2\sqrt{I_1I_2}|\gamma|/(I_1+I_2)$； (4) Cauchy–Schwarz给$|\gamma|\le1$，平衡强度时$V=|\gamma|$。。

\par\medskip

\subsection*{O209\quad 光栅有限缝数主极大宽度}

\textbf{模型与问题：} $N$个等间距$d$的窄缝相干照明，单缝包络在考虑区域近似常数。 (1)建立控制关系；(2)求核心解；(3)分析相关物理量；(4)讨论守恒、极限或稳定性。

\par\textbf{推导要点：} (1) 场振幅为几何级数$A\propto\sin(N\beta)/\sin\beta$，$\beta=\pi d\sin\theta/\lambda$； (2) 主极大在$\beta=m\pi$； (3) 峰附近第一零点间$\Delta\beta=2\pi/N$； (4) 角宽约$\Delta\theta=2\lambda/(Nd\cos\theta_m)$，即随$N$增大而变窄。。

\par\medskip

\subsection*{Q202\quad Landau–Zener避免交叉}

\textbf{模型与问题：} 两能级哈密顿量$H(t)=(vt\sigma_z+\Delta\sigma_x)/2$，$v>0,\Delta>0$，初末时刻趋无穷。 (1)建立控制关系；(2)求核心解；(3)分析相关物理量；(4)讨论守恒、极限或稳定性。

\par\textbf{推导要点：} (1) 瞬时能隙$E_+-E_-=\sqrt{v^2t^2+\Delta^2}$； (2) 最小能隙在$t=0$为$\Delta$； (3) 标准Landau–Zener非绝热跃迁概率$P=\exp[-\pi\Delta^2/(2\hbar v)]$； (4) $v\to0$绝热极限趋零，$v\to\infty$趋1，概率定义需指定绝热或diabatic基底。。

\par\medskip

\subsection*{Q203\quad 量子谐振子相干态}

\textbf{模型与问题：} 升降算符$[a,a^\dagger]=1$，相干态满足$a|\alpha\rangle=\alpha|\alpha\rangle$。 (1)建立控制关系；(2)求核心解；(3)分析相关物理量；(4)讨论守恒、极限或稳定性。

\par\textbf{推导要点：} (1) Fock展开$|\alpha\rangle=e^{-|\alpha|^2/2}\sum_{n\ge0}\alpha^n|n\rangle/\sqrt{n!}$； (2) 粒子数分布$P_n=e^{-|\alpha|^2}|\alpha|^{2n}/n!$； (3) 平均$\bar n=|\alpha|^2$且方差同值； (4) 时间演化$|\alpha(t)\rangle$的幅参数为$\alpha e^{-i\omega t}$（忽略整体相位）。。

\par\medskip

\subsection*{Q204\quad 一维方势阱偶奇束缚态}

\textbf{模型与问题：} 对称有限深势阱$V(x)=-V_0$（$|x|<a$），外部$V=0$，束缚态$-V_0<E<0$。 (1)建立控制关系；(2)求核心解；(3)分析相关物理量；(4)讨论守恒、极限或稳定性。

\par\textbf{推导要点：} (1) 定义$k=\sqrt{2m(E+V_0)}/\hbar$、$\kappa=\sqrt{-2mE}/\hbar$； (2) 偶态连续匹配得$k\tan(ka)=\kappa$； (3) 奇态得$-k\cot(ka)=\kappa$； (4) 阈值由$\kappa\to0$确定额外束缚态出现位置，解须满足$k^2+\kappa^2=2mV_0/\hbar^2$。。

\par\medskip

\subsection*{Q205\quad 双粒子自旋 singlet 相关}

\textbf{模型与问题：} 两个自旋$1/2$粒子处于$|S\rangle=(|\uparrow\downarrow\rangle-|\downarrow\uparrow\rangle)/\sqrt2$。 (1)建立控制关系；(2)求核心解；(3)分析相关物理量；(4)讨论守恒、极限或稳定性。

\par\textbf{推导要点：} (1) 总自旋$(\mathbf S_1+\mathbf S_2)^2|S\rangle=0$； (2) 单粒子约化态均为$I/2$； (3) 任意单位方向$\mathbf a,\mathbf b$的Pauli关联$\langle(\boldsymbol\sigma_1\cdot\mathbf a)(\boldsymbol\sigma_2\cdot\mathbf b)\rangle=-\mathbf a\cdot\mathbf b$； (4) 同向测量结果必反相关，不能由单粒子纯态乘积再现。。

\par\medskip

\subsection*{Q206\quad 方势垒隧穿的透射率}

\textbf{模型与问题：} 一维质量$m$粒子动能$0<E<V_0$穿过宽$a$矩形势垒。 (1)建立控制关系；(2)求核心解；(3)分析相关物理量；(4)讨论守恒、极限或稳定性。

\par\textbf{推导要点：} (1) 势垒外波数$k=\sqrt{2mE}/\hbar$、内衰减$\kappa=\sqrt{2m(V_0-E)}/\hbar$； (2) 波函数及导数在两界面连续； (3) 精确透射率$T=[1+V_0^2\sinh^2(\kappa a)/(4E(V_0-E))]^{-1}$； (4) $\kappa a\gg1$时$T\simeq16E(V_0-E)e^{-2\kappa a}/V_0^2$。。

\par\medskip

\subsection*{Q207\quad Berry相位的自旋几何}

\textbf{模型与问题：} 自旋$1/2$哈密顿量$H=-\frac{\hbar\omega_0}{2}\mathbf n(t)\cdot\boldsymbol\sigma$，$\mathbf n$绝热闭路。 (1)建立控制关系；(2)求核心解；(3)分析相关物理量；(4)讨论守恒、极限或稳定性。

\par\textbf{推导要点：} (1) 瞬时本征态在Bloch球上随$\mathbf n$变化； (2) 动态相位由$-\hbar^{-1}\int E(t)dt$给出； (3) 几何相位为$\gamma_B=-\frac12\Omega_{solid}$（对随场本征态，符号依选择）； (4) 该相位只由闭路包围有向立体角决定，绕行方向反转时变号。。

\par\medskip

\subsection*{Q208\quad 相对论两体非弹碰撞}

\textbf{模型与问题：} 静质量$m$的两个粒子在实验室分别以$+v,-v$对向碰撞，粘成一个复合体。 (1)建立控制关系；(2)求核心解；(3)分析相关物理量；(4)讨论守恒、极限或稳定性。

\par\textbf{推导要点：} (1) 总动量为零，系统质心系即实验室； (2) 总能量$E=2\gamma mc^2$； (3) 复合物静质量$M=E/c^2=2\gamma m$； (4) 额外静能$(M-2m)c^2$来自入射动能，满足四动量守恒。。

\par\medskip

\subsection*{Q209\quad 隧穿双阱有效二能级}

\textbf{模型与问题：} 左右局域态$|L\rangle,|R\rangle$在对称双势阱内的有效哈密顿量$H=E_0I-J\sigma_x$，$J>0$。 (1)建立控制关系；(2)求核心解；(3)分析相关物理量；(4)讨论守恒、极限或稳定性。

\par\textbf{推导要点：} (1) 本征态$(|L\rangle\pm|R\rangle)/\sqrt2$，能量$E_0\mp J$； (2) 初始$|L\rangle$时演化概率$P_R(t)=\sin^2(Jt/\hbar)$； (3) 完全转移时间$\pi\hbar/(2J)$； (4) 引入能量失谐$\Delta\sigma_z/2$时最大转移率下降为$4J^2/(4J^2+\Delta^2)$。。

\par\medskip
